\newpage
\begin{center}
\large\textbf{Минобрнауки России}

\large\textbf{Юго-Западный государственный университет}
\vskip 1em
\normalsize{Кафедра программной инженерии}
\vskip 1em
\ifВКР{
        \begin{flushright}
        \begin{tabular}{p{.4\textwidth}}
        \centrow УТВЕРЖДАЮ: \\
        \centrow Заведующий кафедрой \\
        \hrulefill \\
        \setarstrut{\footnotesize}
        \centrow\footnotesize{(подпись, инициалы, фамилия)}\\
        \restorearstrut
        «\underline{\hspace{1cm}}»
        \underline{\hspace{3cm}}
        20\underline{\hspace{1cm}} г.\\
        \end{tabular}
        \end{flushright}
        }\fi
\end{center}
\vspace{1em}
  \begin{center}
  \large
\ifВКР{
ЗАДАНИЕ НА ВЫПУСКНУЮ КВАЛИФИКАЦИОННУЮ РАБОТУ
  ПО ПРОГРАММЕ БАКАЛАВРИАТА}
  \else
ЗАДАНИЕ НА КУРСОВУЮ РАБОТУ (ПРОЕКТ)
\fi
\normalsize
  \end{center}
\vspace{1em}
{\parindent0pt
  Студента \АвторРод, шифр\ \Шифр, группа \Группа
  
1. Тема «\Тема\ \ТемаВтораяСтрока»
\ifВКР{
утверждена приказом ректора ЮЗГУ от \ДатаПриказа\ № \НомерПриказа
}\fi.

2. Срок предоставления работы к защите \СрокПредоставления

3. Исходные данные для создания программной системы:

3.1. Перечень решаемых задач:}

\renewcommand\labelenumi{\theenumi)}

\begin{enumerate}
	\item Исследовать особенности деятельности редакционно-издательского отдела вуза и процессы проверки учебно-методических материалов.
	\item Проанализировать существующие подходы и программные средства для автоматизации документооборота и проверки документов.
	\item Разработать информационно-аналитическую систему для автоматизации проверки и обработки документов.
	\item Реализовать клиент-серверную архитектуру веб-приложения с использованием REST API.
	\item Реализовать функции авторизации пользователей, загрузки и проверки документов, генерации титульных листов и формирования статистики.	
	\item Провести тестирование разработанной системы и оценить корректность её работы.
\end{enumerate}

{\parindent0pt
	
	3.2. Входные данные и требуемые результаты для программы:}

\begin{enumerate}
	\item Входными данными для программной системы являются документы формата DOCX, данные пользователей, параметры проверки документов и информация, вводимая через веб-интерфейс системы.
	
	\item Выходными данными являются результаты проверки документов, PDF-отчёты об ошибках, DOCX-файлы с выделенными нарушениями, сгенерированные титульные листы, статистическая информация и данные, отображаемые в пользовательском интерфейсе.
\end{enumerate}

{\parindent0pt
	
	4. Содержание работы (по разделам):
	
	4.1. Введение.
	
	4.2. Анализ предметной области.
	
	4.3. Техническое задание: основание для разработки, назначение разработки, требования к программной системе, требования к оформлению документации.
	
	4.4. Технический проект: общие сведения о программной системе, проектирование архитектуры клиент-серверной системы, проектирование пользовательского интерфейса, описание взаимодействия клиентской части с REST API.
	
	4.5. Рабочий проект: спецификация компонентов и классов программной системы, тестирование программной системы, сборка компонентов программной системы.
	
	4.6. Заключение.
	
	4.7. Список использованных источников.

5. Перечень графического материала:

\списокПлакатов

\vskip 2em
\begin{tabular}{p{6.8cm}C{3.8cm}C{4.8cm}}
Руководитель \ifВКР{ВКР}\else работы (проекта) \fi & \lhrulefill{\fill} & \fillcenter\Руководитель\\
\setarstrut{\footnotesize}
& \footnotesize{(подпись, дата)} & \footnotesize{(инициалы, фамилия)}\\
\restorearstrut
Задание принял к исполнению & \lhrulefill{\fill} & \fillcenter\Автор\\
\setarstrut{\footnotesize}
& \footnotesize{(подпись, дата)} & \footnotesize{(инициалы, фамилия)}\\
\restorearstrut
\end{tabular}
}

\renewcommand\labelenumi{\theenumi.}
